% !TeX root = ../main.tex

% ================ 結論與未來展望 ================= %
\section{結論與未來展望}

% --------------------- 結論 --------------------- %
\subsection{結論}
面對全球淨零碳排之迫切趨勢，5G-Advanced 與未來 6G 網路在追求極致傳輸速率的同時，亦面臨著大規模天線陣列（Massive MIMO）帶來的巨大功耗挑戰。傳統之空間域適應（SDA）技術多仰賴歷史負載統計進行反應式天線降維，在面對高度非平穩之突發流量時，往往陷入決策滯後、頻繁切換之乒乓效應，以及排隊延遲發散等效能瓶頸。為突破此一困境，本研究立基於 3GPP TR 38.864 國際標準之功耗框架，創新提出並實作了一套流量感知動態天線調適機制。

回顧本論文之研究歷程，從底層數學模型之建構、邏輯沙盒之參數調校，至最終 C++ 系統級模擬器（SLS）之全網實戰驗證，本研究之核心學術貢獻與工程實證成果可歸納為以下四大維度：

\begin{enumerate}
    \item \textbf{導入並改良李雅普諾夫最佳化理論，建立嚴謹之 QoS 數學邊界：} 
    本研究成功將 5G 基地台之 MAC 層排程與 PHY 層天線狀態，轉化為具備數學收斂保證之隨機網路最佳化問題。透過漂移加懲罰理論，演算法僅需依賴當前之瞬時佇列積壓與通道狀態即可進行線上決策，徹底擺脫了傳統歷史視窗之滯後性。更重要的是，本機制在理論上嚴格保證了系統能以 $\mathcal{O}(1/V)$ 之速率逼近絕對節能極限，同時確保最壞情況下的佇列延遲被限制在 $\mathcal{O}(V)$ 的線性安全上界內，從根本上杜絕了網路癱瘓之風險。    
    \item \textbf{首創動態 $V$ 值流量感知機制與非線性參數調優：} 
    為克服傳統靜態權重在突發流量下缺乏彈性之缺陷，本研究於決策引擎中引入了「容量預警餘裕（Capacity Margin）」與「非線性形狀因子（Shape Factor）」，設計出具備流量感知能力之動態 $V$ 值自適應機制。透過建構基於真實 5G 串流影音（Stored Streaming, 依據 Netflix 實際網路軌跡）驅動之 MATLAB 演算法沙盒，本研究精確量化了不同參數在能效與延遲之間的最佳化折衷邊界。實驗證實，設定非線性調節因子 $\gamma = 2.0$，能使系統在「輕載深度休眠」與「重載瞬間喚醒」之間取得最佳之動態折衷，確保基地台面對各類極端負載皆能游刃有餘。    
    \item \textbf{3GPP 大規模系統級模擬 (SLS) 之極致節能實證：} 
    本研究克服了系統底層架構之高度複雜性，將最佳化模組完整移殖至具備 7 Sites、21 Cells 之 Wrap-around 網格拓樸 SLS 模擬平台。在包含 3GPP UMa 三維通道衰落模型（3D Spatial Channel Model）與多細胞同頻干擾的惡劣環境下，與傳統負載導向機制及靜態演算法進行了深度交叉對比。量化數據強而有力地證實，本研究所提之動態機制在低負載與輕負載情境下，斬獲了高達 30.3\% 之實質節能率（ESG），展現了傲視現有方案之冗餘能耗削減能力。    
    \item \textbf{跨層實體干擾與排程紅利之深層物理機制剖析：} 
    本研究並未止步於實驗數據之表象，更透過深度之累積分布函數剖析，揭露了演算法於現網環境下運作之底層通訊原理。首先，本研究精準指出了實體層波束拓寬外溢與媒體存取控制層通道狀態資訊過期估測之間，交織引發之資訊不對稱性，客觀合理解釋了演算法在追求極致節能時所付出之用戶感知吞吐量代價。其次，透過跨負載情境之對比，驗證了「多使用者多樣性增益」與李雅普諾夫決策引擎中\textbf{「基於佇列積壓驅動之全天線組態重構機制」}的聯合效應，從而定量且科學地解釋了輕負載情境下反而能有效消弭長尾延遲之自洽現象。這些物理維度之觀察與剖析，為未來跨層資源調配演算法之設計提供了無可取代的學術與工程參考價值。
\end{enumerate}

綜合上述成果，本研究提出之機制成功在「數學理論保證」與「工業實務部署」之間搭建了橋樑。其卓越的省電能力與「主動容忍可控延遲、精準拒絕效能溢出」的獨特演算法性格，使其極度契合未來 5G-Advanced 中針對大規模機器聯網（mMTC）、智慧物聯網，以及對營運電費與碳足跡高度敏感之「綠能背景網路切片（Green Network Slicing）」應用。本論文之量化實證，為未來永續通訊系統（Sustainable Communication Systems）之底層架構設計，奠定了極具前瞻性與工程落地價值之基石。

% ------------------- 未來展望 ------------------- %
\subsection{未來展望}
基於本研究於系統級模擬中所獲取之實驗數據與物理機制剖析，針對基地台之空間域適應與電力能效最佳化技術，未來研究可朝以下幾個具備工程可行性之維度進行深化與延伸：

\begin{enumerate}
    \item \textbf{環境感知之非線性調節因子自適應優化：} 本研究已實證，固定之非線性形狀因子 $\gamma = 2.0$ 雖能於沙盒中取得優異之統計折衷，然而在多細胞強干擾環境下，若遭遇相鄰細胞同步採取省電策略，極易因決策資訊過期而引發過度降階。未來研究可嘗試將固定常數 $\gamma$ 升級為「具備環境干擾感知能力之動態調節器」。當本地基地台偵測到相鄰細胞負載急遽上升、或是本地訊號干擾雜訊比（SINR）發生劇烈震盪時，演算法能即時動態調降 $\gamma$ 之數值，藉此強迫系統回歸保守之高階天線組態（如 64T），以維持窄波束之空間抗干擾能力，在真實佈建中更穩健地保障連線服務品質。

    \item \textbf{結合深度強化學習之動態控制參數預測：} 李雅普諾夫架構之核心優勢在於無需掌握流量之先驗知識即可執行毫秒級之線上即時決策，然而其在本質上缺乏對中長期趨勢（如日夜間通勤之流量潮汐變化）之巨觀記憶與預測能力。未來研究可探討將深度強化學習（Deep Reinforcement Learning, DRL）或長短期記憶網路（LSTM）等模型與李雅普諾夫引擎進行混合架構設計（Hybrid Control Architecture）。由 AI 模型負責在較長之時間尺度下，預測未來時間視窗之流量基底並輸出前置之演算法邊界參數，再由李雅普諾夫引擎於短時間尺度（TTI 級別）執行瞬時之干擾迴避與天線維度微調，以此雙層控制架構進一步突破永續網路之節能極限。

    \item \textbf{引入鄰區長期平均負載歷史之悲觀干擾補償折扣：}
    在現有分散式決策架構下，單一基地台難以即時獲知鄰近細胞的天線切換狀態。為彌補此一資訊盲區，未來可於預期傳輸速率估測模型中，引入一基於鄰區長期平均負載歷史的悲觀干擾補償折扣參數（$\delta_{\text{leak}}$）。該參數將作為本地決策的數學防護網，使調整後之預期速率可定義為：
    \begin{equation}
    R_{n, \text{adjusted}}(a(t)) = R_n(a(t)) \cdot (1 - \delta_{\text{leak}}(a(t), \bar{\rho}_{\text{neighbor}}))
    \end{equation}
    其中 $\delta_{\text{leak}}$ 的大小應與天線組態之波束寬度及鄰區歷史平均負載 $\bar{\rho}_{\text{neighbor}}$ 成正比（例如，16T 寬波束配置下的 $\delta_{\text{leak}}(16\text{T})$ 應大於 64T 窄波束的 $\delta_{\text{leak}}(64\text{T})$）。藉由在李雅普諾夫決策引擎中預先扣除波束拓寬可能引發的干擾外溢代價，可有效修正單一細胞通道狀態資訊（CSI）估測的樂觀偏差，提升全網動態干擾博弈的穩健性。

    \item \textbf{基於 O-RAN RIC 架構之準即時多細胞區域聯合優化機制：}
    考量多細胞間若要在毫秒（ms）級別做出完全即時的協同反應，在工程佈建上本來就有回程延遲（Backhaul Latency）而不切實際的物理瓶頸。為突破此侷限，未來研究可結合開放式無線接取網路（O-RAN）架構，將李雅普諾夫優化引擎部署於近即時無線接取網路智慧控制器（Near-RT RIC）中，作為獨立的擴充應用程式（xApp）。透過標準化之 E2 介面，週期性地聚合區域內多細胞的佇列積壓與干擾矩陣資訊，將原本分散式的決策盲區，轉換為準即時（Near-RT）的區域聯合優化機制，藉此在可容忍的信令延遲下，實現全網能效與干擾抑制的全局最佳化。
\end{enumerate}